% !TEX root = ../../tree_notes.tex
% =====================================================================
%  5.3  Path-as-Sequence and the Depth Sequence             sec:c24-path
%  Exports: def:c24-path  prop:c24-sweep
%           lst:c24-path-stack  ex:c24-path-run
%           prob:c24-levelsum  lst:c24-levelsum
%           prob:c24-kthlevel  lst:c24-kthlevel  practice: LC 1367 (unlabelled)
% =====================================================================
\section{Path-as-Sequence and the Depth Sequence}\label{sec:c24-path}

\begin{previouslybox}
\begin{enumerate}[nosep, leftmargin=1.6em]
  \item Which nodes of $G$ share a depth, and what is $\depth$ for every node? (\figref{c04-running-g})
  \item In the DFS, when you \emph{enter} and \emph{leave} a node, what does the recursion
        stack hold at that instant, in terms of the root-to-$v$ path? (\secref{c24-times})
  \item State the subtree range $[\tin(v), \tout(v))$; why is it the \emph{opposite} end of
        the same stick from a root-path? (\secref{c24-flatten})
  \item Why did we say the 2-tick (entry \emph{and} exit) clock is the one that unlocks path
        sums? (\secref{c24-times}, variation 3)
\end{enumerate}
\end{previouslybox}

\begin{topicintro}
Flattening answered \emph{subtrees}; this section answers \emph{paths}. Read the same DFS one
way and a subtree is a contiguous \emph{block}; read it another way and the recursion stack
\emph{is} the root-to-current path, a sequence you can push onto and pop off for free. Two
views fall out. (i) The \term{root-path view}: carry a stack (or a $\pm w$ difference array on
the Euler clock) so that at any node you can read the aggregate along its path from the root
in $O(1)$. (ii) The \term{depth sequence}: bucket nodes by $\depth$ and process level by
level -- the array behind every ``per-level'' problem. LC 1161 and LC 2583 are exactly this:
sum each level, then take the best or the $k$-th largest.
\textbf{Prerequisites.} \secref{c24-times} ($\tin/\tout$, both clocks); \secref{c24-flatten}
(the subtree range, for contrast); \chref{c03} BFS levels.

\textbf{After this section you can}
\begin{itemize}[nosep]
  \item maintain the root-to-node path as a DFS stack and answer path aggregates; \hfill\lv{3}
  \item turn ``add $w[v]$ to a whole subtree, then read the root-path to $u$'' into an
        in/out difference sweep;                                          \hfill\lv{4}
  \item build a \emph{level array} (buckets by $\depth$) and aggregate per level;   \hfill\lv{2}--\lv{3}
  \item find the best / $k$-th level sum by combining a BFS with a sort.           \hfill\lv{3}
\end{itemize}
\end{topicintro}

\begin{recognitionbox}[Pattern card -- paths and levels]
\reckey{the question is about the \emph{root-to-node path} (``sum of values on the way down'',
``is there a root-path matching this pattern''), or about a \emph{level} (``the level with the
largest sum'', ``average per level'', ``$k$-th largest level sum'').}
\recuse{
  \item ``path from the root to $v$'' aggregates asked for many $v$ -- carry a stack or a
        difference array (\propref{c24-sweep});
  \item ``per-level sum / max / average'', ``which level is heaviest'' -- a level array
        bucketed by $\depth$ (LC 1161, \probref{c24-levelsum});
  \item ``$k$-th something over the levels'' -- collect level sums, select (LC 2583,
        \probref{c24-kthlevel});
  \item ``does this \emph{sequence} appear going downward from some node'' -- DFS carrying the
        pattern along the path stack (the LC 1367 practice below).
}
\recnotuse{
  \item aggregating over a \emph{subtree} -- that is \secref{c24-flatten}, not here;
  \item the path between \emph{two arbitrary} nodes (root-path $\to$ LCA $\to$ distance) --
        \secref{c24-dist} and \chref{c11};
  \item path \emph{with updates} as a first-class structure -- heavy-light (\chref{c16}).
}
One-line test: \emph{is the unit of work ``down from the root'' (a path) or ``across one
depth'' (a level)? Then use the path stack or the level array -- this section.}
\end{recognitionbox}

\subsection{Intuition}

The recursion stack of a DFS is, at every instant, exactly the path from the root to the node
you are standing on: pushing a node extends the path, popping it retracts. So the ``sequence''
you carry down (\secref{c02-topdown}) is not extra bookkeeping -- it \emph{is} the stack.
Everything you want to know about the root-path (its sum, whether it contains a pattern, its
longest increasing run) can be maintained incrementally as you descend and undone as you
ascend.

\begin{figure}[H]
\centering
\begin{tikzpicture}[diagram]
  \node[troot] (n0) at (0,0) {$0$};
  \node[tnode] (n1) at (-3,-1.4) {$1$};
  \node[tnode] (n2) at (0,-1.4) {$2$};
  \node[tnode] (n3) at (3,-1.4) {$3$};
  \node[tleaf] (n4) at (-3.9,-2.8) {$4$};
  \node[tleaf] (n5) at (-2.1,-2.8) {$5$};
  \node[tleaf] (n6) at (3,-2.8) {$6$};
  \draw[tedge] (n0) -- (n1); \draw[tedge] (n0) -- (n2); \draw[tedge] (n0) -- (n3);
  \draw[tedge] (n1) -- (n4); \draw[tedge] (n1) -- (n5); \draw[tedge] (n3) -- (n6);
  \draw[tedgehl, line width=1.8pt] (n0) -- (n1) -- (n4);
  \node[annot, align=left, font=\scriptsize] at (-6.5,-0.4)
    {standing at $4$, the\\stack $=[0,1,4]$\\$=\depth$-path from root};
  \node[annot, align=left, font=\scriptsize] at (4.3,-2.3)
    {levels: $L_0=\{0\}$\\$L_1=\{1,2,3\}$\\$L_2=\{4,5,6\}$};
\end{tikzpicture}
\caption{Two readings of one DFS. The \textbf{highlighted chain} $0,1,4$ is the recursion
stack at node $4$ -- the root-path, a \emph{sequence}. The \textbf{horizontal bands} are the
level arrays $L_d=\{v:\depth(v)=d\}$ -- the \emph{depth sequence}. Paths read down a column;
levels read across a row.}\label{fig:c24-path}
\end{figure}

Now bucket nodes by depth. A BFS naturally produces these \term{level bands}; a DFS can push
each node into \code{level[$\depth(v)$]}. Once you have the array of per-level sums (or
counts, or maxima) you can scan it for the best level (\probref{c24-levelsum}), sort it for
order statistics (\probref{c24-kthlevel}), or average it -- all offline from the tree, on a
plain array.

\subsection{Definitions and Notation}

\begin{definition}{Root-path; level array}{c24-path}
The \term{root-path} of $v$ is the sequence of nodes from the root to $v$ inclusive, which a
DFS holds as its call stack when it reaches $v$. The \term{level} (or \term{depth bucket})
$L_d = \{\, v : \depth(v) = d \,\}$; the \term{depth sequence} is the partition
$(L_0, L_1, \dots, L_{\height})$. A \emph{path aggregate} $A(v)$ (sum, xor, list of values) is
a quantity defined over the root-path of $v$.
\end{definition}

\noindent\textbf{In plain words.} ``Path from the root to $v$'' $=$ the stack you had to build
to reach $v$; maintain it as you go. ``Level $d$'' $=$ everyone exactly $d$ edges below the
root; group them and run one aggregate across each group.

\begin{proposition}{Root-path sums by an in/out difference sweep}{c24-sweep}
Give each node a weight $w[v]$ and use the \emph{both-tick} Euler clock of \secref{c24-times}
(values in $[0, 2n)$). Build a difference array $D$ of length $2n$ with $D[\tin(v)] += w[v]$
and $D[\tout(v)] -= w[v]$. Then the prefix sum $\sum_{i \le \tin(u)} D[i]$ equals
$\sum_{x} w[x]$ over the \emph{root-path to $u$}. Hence path sums -- and, dually, ``add to a
subtree then read a path'' questions -- reduce to a prefix scan over the tour.
\end{proposition}
\begin{proofbox}
$w[x]$ is $+1$-weighted at $\tin(x)$ and cancelled at $\tout(x)$. The prefix sum at position
$\tin(u)$ therefore counts $w[x]$ iff $\tin(x) \le \tin(u) < \tout(x)$, i.e.\ iff $x$'s
interval contains the instant we enter $u$. By \thmref{c24-anc} those $x$ are exactly the
ancestors of $u$ (including $u$) -- the root-path. Linearity gives the weighted sum. The
subtree-add / point-query case is the mirror image: add over $[\tin,\tout)$ and read at
$\tin(u)$, which the same $D$ supports with the roles swapped.
\end{proofbox}
\begin{insightbox}
This is the \emph{dual} of \secref{c24-flatten}. There, a subtree was a range and a node was a
point; here, a root-path is a prefix of the tour and a node is again a point. ``Subtree
update $+$ path query'' and ``path update $+$ subtree query'' are the two compositions these
two sections give you -- the full generality (with both operations online) is a heavy-light
or Euler-tour-tree affair (\chref{c16}, \chref{c23}).
\end{insightbox}

\subsection{The Algorithm: carry the path, or bucket the levels}

\begin{lstlisting}[caption={DFS carrying a root-path stack, and a DFS/BFS filling a level array}, label={lst:c24-path-stack}]
#include <bits/stdc++.h>
using namespace std;

vector<vector<int>> adj;
vector<vector<int>> levelVec;   // levelVec[d] = node ids at depth d
vector<int> depth_;

// (A) per-level buckets via DFS (push node id into its depth's list)
void dfsLevels(int u, int p, int d) {
    if ((int)levelVec.size() <= d) levelVec.resize(d + 1);
    levelVec[d].push_back(u);
    depth_[u] = d;
    for (int v : adj[u]) if (v != p) dfsLevels(v, u, d + 1);   // parent guard
}

// (B) root-path stack: aggregate along the way down, undone on the way up
long long pathSum = 0, bestAtNode_toRoot = LLONG_MIN;
vector<long long> pathVals;      // the current root-to-node sequence
void dfsPath(int u, int p, long long w) {
    pathSum += w;                // ENTER: extend the path aggregate
    pathVals.push_back(u);       //          and the path sequence
    bestAtNode_toRoot = max(bestAtNode_toRoot, pathSum);   // answer for THIS node
    for (int v : adj[u]) if (v != p) dfsPath(v, u, w);
    pathSum -= w;                // EXIT: retract exactly what we added
    pathVals.pop_back();
}
\end{lstlisting}
\begin{complexitybox}
Both passes are $\bigO{n}$ (each node entered/left once); the level array and depth array
are $\bigO{n}$ space total. The path stack costs $\bigO{h}$ live entries. Answering ``sum on
root-path to $u$'' is then $O(1)$ if you recorded it during the descent, or $\bigO{\depth(u)}$
if you replay \code{pathVals}. Recursion is $\bigO{h}$ deep -- iterative (explicit stack with
a matching value stack) on a chain.
\end{complexitybox}

\subsection{Worked Example}

\begin{example}{Levels and root-paths of $G$ (with $w[v]=v$)}{c24-path-run}
$\depth=[0,1,1,1,2,2,2]$, so the level buckets are
$L_0=\{0\}$, $L_1=\{1,2,3\}$, $L_2=\{4,5,6\}$, giving per-level sums
$S_0=0$, $S_1=1{+}2{+}3=6$, $S_2=4{+}5{+}6=15$. The heaviest level is $L_2$
($1$-indexed ``level $3$''), and the sorted level sums $\{15,6,0\}$ make the $2$-nd largest
$6$. Now the \emph{root-path} to node $4$: the stack when DFS stands at $4$ is
$[0,1,4]$, path sum $0{+}1{+}4=5$; to node $6$: stack $[0,3,6]$, sum $9$. The in/out sweep of
\propref{c24-sweep} produces the same $5$ and $9$ as prefix sums at $\tin(4)$ and $\tin(6)$.
\end{example}
\begin{tracebox}[Push/pop of the path stack along the DFS order $0,1,4,5,2,3,6$]
Enter $0$ (path $[0]$, sum $0$); enter $1$ ($[0,1]$, sum $1$); enter $4$ ($[0,1,4]$, sum $5$)
\textbf{answer@4}$=5$; leave $4$; enter $5$ ($[0,1,5]$, sum $6$) \textbf{answer@5}$=6$; leave
$5$; leave $1$; enter $2$ ($[0,2]$, sum $2$) answer@2$=2$; leave $2$; enter $3$ ($[0,3]$,
sum $3$); enter $6$ ($[0,3,6]$, sum $9$) answer@6$=9$; leave $6$; leave $3$; leave $0$. Every
``answer@node'' is the root-path sum -- read at the instant of the push. \checkmark\
\figref{c24-path}.
\end{tracebox}

\subsection{Variations and Twists}

\begin{itemize}[nosep]
  \item \textbf{Per-level max / count / average}: swap the $+$ in the bucket fold for
        $\max$, $1$, or keep sum and count for an average (\chref{c03}'s level-order family). \hfill\lv{2}
  \item \textbf{Path with a predicate} (increasing run, distinct values): carry the
        \code{pathVals} stack and maintain the statistic incrementally on push/pop. \hfill\lv{3}--\lv{4}
  \item \textbf{Subtree-add then root-path-read}: the in/out difference sweep,
        \propref{c24-sweep}, is exactly the two-composition. \hfill\lv{4}
  \item \textbf{Any-pair path}: not from the root -- get LCA and stitch two root-paths;
        \secref{c24-dist} and \chref{c11}.                                   \hfill\lv{4}--\lv{5}
  \item \textbf{Level-order without a queue}: fill $\depth$ by DFS, then sort nodes by
        $\depth$ -- the ``depth sequence'' is a lazy BFS.                     \hfill\lv{3}
\end{itemize}

\subsection{Common Pitfalls}\label{sec:c24-path-pitfalls}

\begin{warningbox}[Trap 1 -- forgetting to undo on the way up]
A path aggregate pushed on entry must be \emph{popped} on exit, or sibling subtrees inherit a
phantom prefix and every answer below is inflated.
\end{warningbox}
\begin{lstlisting}[caption={\textcolor{warningred}{WRONG}: push on entry, never retract}, language=C++, basicstyle=\ttfamily\small]
void dfs(int u,int p){ pathSum+=w[u]; for(int v:adj[u])if(v!=p)dfs(v,u); /* no -= */ }
\end{lstlisting}
\begin{lstlisting}[caption={\textcolor{successgreen}{RIGHT}: subtract exactly what you added}, language=C++, basicstyle=\ttfamily\small]
void dfs(int u,int p){ pathSum+=w[u];
    for(int v:adj[u])if(v!=p)dfs(v,u);
    pathSum-=w[u]; }                     // restore for the sibling branch
\end{lstlisting}

\begin{warningbox}[Trap 2 -- off-by-one indexing levels]
Root is $\depth=0$ but ``level'' in problems is usually $1$-indexed. Reporting
$\arg\max$ of the level array without $+1$ is a classic wrong answer (LC 1161).
\end{warningbox}
\begin{lstlisting}[caption={\textcolor{warningred}{WRONG}: returns a $0$-based depth as the level}, language=C++, basicstyle=\ttfamily\small]
return bestDepth;                 // problem wants level = bestDepth + 1
\end{lstlisting}
\begin{lstlisting}[caption={\textcolor{successgreen}{RIGHT}: shift to the requested indexing}, language=C++, basicstyle=\ttfamily\small]
return bestDepth + 1;             // root at depth 0 is "level 1"
\end{lstlisting}

\begin{warningbox}[Trap 3 -- resizing / touching \code{level[d]} out of range]
On the first visit to a new maximum depth the bucket does not exist; index it only after
resizing, or size the array to $\height{+}1$ up front. A stale bound is an out-of-range read.
\end{warningbox}
\begin{lstlisting}[caption={\textcolor{warningred}{WRONG}: push into a bucket that may not exist}, language=C++, basicstyle=\ttfamily\small]
levelVec[d].push_back(u);                 // d may exceed current size -> UB
\end{lstlisting}
\begin{lstlisting}[caption={\textcolor{successgreen}{RIGHT}: grow the vector to depth first}, language=C++, basicstyle=\ttfamily\small]
if ((int)levelVec.size() <= d) levelVec.resize(d + 1);
levelVec[d].push_back(u);
\end{lstlisting}

\subsection{Problem Lab}

Two level-array problems, ascending. LC 1161 \emph{Maximum Level Sum} is the warm-up: fill the
depth buckets, take the heaviest, return its (1-indexed) level. LC 2583 \emph{Kth Largest Sum
in a Binary Tree} adds an order statistic on the same array. Both are the depth sequence
$\longrightarrow$ an ordinary array $\longrightarrow$ one scan or sort.

% ---------------------------------------------------------------------------
\solvedproblem{c24-levelsum}{Maximum Level Sum of a Binary Tree}{LeetCode 1161}
  {https://leetcode.com/problems/maximum-level-sum-of-a-binary-tree/}{LC Medium L2}
  {BFS/DFS level buckets, per-level sum, argmax with 1-indexing}

\begin{problemstatement}[Maximum Level Sum of a Binary Tree]
Given the \code{root} of a binary tree, the \textbf{level} of a node is its distance from the
root with the root at level $1$. The \term{sum} of a level is the sum of node \emph{values} at
that level. Return the \emph{smallest} level $X$ whose level-sum is maximal. Node values can
be \textbf{negative}. $1 \le$ number of nodes $\le 10^4$.
\end{problemstatement}

\solstep{Understanding the Problem}
This is the depth sequence made concrete: bucket every node by its level, sum each bucket, and
report the earliest level attaining the maximum. Because node values may be negative, the
maximum level-sum is \emph{not} necessarily the root level -- you must actually compute every
level's sum, and initialise the running best to $-\infty$ (not $0$).

\solstep{Key Observations}
\begin{observation}
A BFS drains one level at a time, so the ``current queue size'' is the level's population --
sum those, then move on. Alternatively a DFS pushes each value into \code{level[d]}. Either
yields the depth sequence.
\end{observation}
\begin{observation}
Ties go to the \emph{smallest} level, so scan levels in increasing order and only update the
best on a \textbf{strictly} larger sum (use $>$, not $\ge$).
\end{observation}
\begin{observation}
Negative values make a $0$-initialised best wrong on the very first level; seed with the
first level's sum (or \code{LLONG\_MIN}) and track the level index alongside.
\end{observation}

\solstep{Solution}
BFS by levels: for each level, drain exactly the current queue length, accumulating a sum; if
that sum exceeds the running best, record this level number (1-indexed). Return the recorded
level.

\solstep{Proof of Correctness}
\begin{proofbox}
Level-order BFS visits nodes grouped by depth in increasing depth (the queue discipline of
\chref{c03} 3.3), and draining the front of size $|L_d|$ per round enumerates exactly $L_d$,
so the loop computes each level's true sum and no node is double- or mis-counted. Scanning in
increasing level and updating only on a \emph{strict} improvement keeps the smallest level
attaining the maximum, matching the problem's tie-break. Since the best is seeded below every
reachable sum, the maximum over all levels (possibly negative) is found. The returned index is
$1$-indexed as required (Trap 2).
\end{proofbox}

\solstep{Implementation}
\begin{lstlisting}[caption={Maximum Level Sum via level BFS (LeetCode 1161)}, label={lst:c24-levelsum}]
#include <bits/stdc++.h>
using namespace std;

struct TreeNode { int val; TreeNode *left, *right;
    TreeNode(int v):val(v),left(nullptr),right(nullptr){} };

class Solution {
public:
    int maxLevelSum(TreeNode* root) {
        long long best = LLONG_MIN; int bestLevel = 1, level = 0;
        queue<TreeNode*> q; q.push(root);
        while (!q.empty()) {
            ++level;
            long long sum = 0;
            int cnt = (int)q.size();           // exactly this level's width
            for (int i = 0; i < cnt; ++i) {    // drain one level
                TreeNode* u = q.front(); q.pop();
                sum += u->val;
                if (u->left)  q.push(u->left);
                if (u->right) q.push(u->right);
            }
            if (sum > best) { best = sum; bestLevel = level; }   // strict: smallest wins
        }
        return bestLevel;                      // 1-indexed level
    }
};
\end{lstlisting}

\solstep{Trace}
\begin{tracebox}[Binary running tree $R$ (\figref{c03-running}), values $1..6$]
Level 1: $\{1\}$ sum $1$, best $1$/lvl $1$. Level 2: $\{2,3\}$ sum $5$, best $5$/lvl $2$.
Level 3: $\{4,5,6\}$ sum $15$, best $15$/lvl $3$. Return $\mathbf{3}$. \checkmark\
\probref{c24-levelsum}.
\end{tracebox}

\solstep{Complexity}
\begin{complexitybox}
Time $\bigO{n}$ (each node dequeued once), space $\bigO{w}$ the widest level ($\le n$). A
DFS variant with a \code{level[]} array is also $\bigO{n}$ but $\bigO{h}$ stack. Verdict: \ac.
\end{complexitybox}

\solstep{Pitfalls and Variants}
\begin{warningbox}
Initialising \code{best=0} fails when all levels are negative; use \code{LLONG\_MIN}. Using
$\ge$ instead of $>$ returns the \emph{largest} level among ties, not the smallest. And
remember level is $1$-indexed here while BFS/DFS depth is often $0$-indexed.
\end{warningbox}

\begin{insightbox}[Key takeaways]
\begin{enumerate}[nosep, leftmargin=1.6em]
  \item ``Per level'' $=$ the depth sequence $=$ one array bucketed by $\depth$; the tree
        disappears once you have it.
  \item Draining \code{q.size()} per round is the clean way to get level widths.
  \item Tie-break and indexing (smallest level, $1$-based) are where the WA lives, not the sum.
\end{enumerate}
\end{insightbox}

% ---------------------------------------------------------------------------
\solvedproblem{c24-kthlevel}{Kth Largest Sum in a Binary Tree}{LeetCode 2583}
  {https://leetcode.com/problems/kth-largest-sum-in-a-binary-tree/}{LC Medium L3}
  {level sums + nth\_element/sort, not-enough-levels edge}

\begin{problemstatement}[Kth Largest Sum in a Binary Tree]
Given the \code{root} of a binary tree and an integer $k$. The \term{level sum} is the sum of
node \emph{values} at a level (root at level $1$). Return the $k$-th \textbf{largest} level sum
(nozzling that level sums are counted with multiplicity). If the tree has fewer than $k$
levels, return $-1$. Node values can be negative. $1 \le n \le 10^5$, $1 \le k \le 10^5$.
\end{problemstatement}

\begin{attemptgate}[Attempt it yourself first -- \lv{3}, give yourself 20 minutes]
A fair first attempt: reuse the level-sum machinery of \probref{c24-levelsum}, then ask how to
pull out the $k$-th largest of the resulting array.
\peek{you now have one number per level -- how many levels could there be, versus $n$?}
\nudge{collect all level sums into a vector, then either sort descending or use
\code{nth\_element}. When is the answer $-1$?}
\fullhint{BFS by levels into a \code{vector<long long> sums}; if \code{sums.size()}$<k$
return $-1$; else sort descending and return \code{sums[k-1]} (or \code{nth\_element} for
$\bigO{\text{levels}}$). Values are \code{long long} (level sums reach $\sim n\cdot10^5$ with
negatives). Full build below.}
\end{attemptgate}

\solstep{Understanding the Problem}
Identical front end to \probref{c24-levelsum} -- build the depth sequence and sum each level.
The only new question is an \emph{order statistic} on the array of level sums: the $k$-th
largest. Because there are at most $\height+1 \le n$ levels, the array is small and a sort is
trivially fast; the edge case is ``fewer than $k$ levels $\Rightarrow -1$''.

\solstep{Key Observations}
\begin{observation}
Number of levels $\le n$ (and $\le$ nodes), so sorting all level sums is at worst
$\bigO{n\log n}$ and typically $\bigO{\height\log\height}$ -- negligible next to the BFS.
\end{observation}
\begin{observation}
Level sums can be negative and can reach $\pm\,n\cdot 10^{5}$, so store in \code{long long}.
\end{observation}
\begin{observation}
``$k$-th largest'' with $k$ up to $10^5$ and possibly $k > $ number of levels: guard the size
before indexing.
\end{observation}

\solstep{Think It Through With Me}
\begin{thinkbox}
\deadend{Re-run a BFS $k$ times hunting the $i$-th largest each pass.}{There is only \emph{one}
array of level sums; compute it once and select from it. Re-traversing the tree is $\bigO{kn}$
for nothing.}
\deadend{Keep a running min-heap of size $k$ during the BFS levels.}{Correct and $\bigO{n +
\height\log k}$, but pointless complexity -- the level count is tiny; just sort. The heap
shines when the stream is huge and you cannot store it, which is not the case here.}
\aha{Collect all level sums into a \code{vector<long long>}, sort descending, return
\code{[k-1]} if it exists, else $-1$.}
\end{thinkbox}

\solstep{The Approach Ladder}
\begin{center}\renewcommand{\arraystretch}{1.3}
\begin{tabular}{@{}c p{0.235\linewidth} p{0.28\linewidth} l c@{}}
\toprule
\textbf{\#} & \textbf{Approach} & \textbf{Idea} & \textbf{Time} & \textbf{Verdict} \\
\midrule
1 & \lv{3} level sums + sort desc & one BFS, then sort & $O(n+\ell\log\ell)$ & \ac \\
2 & \lv{4} level sums + nth\_element & linear select & $O(n+\ell)$ & \ac \\
3 & \lv{2} BFS $k$ times & re-scan tree per rank & $O(kn)$ & \tle \\
\bottomrule
\end{tabular}\end{center}
$\ell$ is the level count ($\le n$, usually far smaller). Rungs 1 and 2 both pass; 1 is the
readable default.

\solstep{Solution}
BFS by levels filling \code{sums}. If \code{sums.size()}$<k$ return $-1$. Otherwise sort
\code{sums} in descending order and return \code{sums[k-1]}.

\solstep{Proof of Correctness}
\begin{proofbox}
By the argument in \probref{c24-levelsum}'s proof, the level BFS produces exactly one sum per
level $\{S_1,\dots,S_\ell\}$ (root at level $1$), each equal to the true level sum. Sorting
these $\ell$ values in non-increasing order places the $k$-th largest at index $k-1$ when
$k \le \ell$; when $k > \ell$ there is no $k$-th level and the problem mandates $-1$, which
the size guard returns. Hence the output matches the definition in both branches.
\end{proofbox}

\solstep{Implementation}
\begin{lstlisting}[caption={Kth Largest Level Sum (LeetCode 2583)}, label={lst:c24-kthlevel}]
#include <bits/stdc++.h>
using namespace std;

struct TreeNode { long long val; TreeNode *left, *right;
    TreeNode(long long v):val(v),left(nullptr),right(nullptr){} };

class Solution {
public:
    long long kthLargestLevelSum(TreeNode* root, int k) {
        vector<long long> sums;
        queue<TreeNode*> q; q.push(root);
        while (!q.empty()) {
            long long s = 0;
            int cnt = (int)q.size();
            for (int i = 0; i < cnt; ++i) {           // drain one level
                TreeNode* u = q.front(); q.pop();
                s += u->val;
                if (u->left)  q.push(u->left);
                if (u->right) q.push(u->right);
            }
            sums.push_back(s);
        }
        if ((int)sums.size() < k) return -1;          // fewer than k levels
        nth_element(sums.begin(), sums.begin() + (k - 1), sums.end(),
                    greater<long long>());            // k-th largest into place
        return sums[k - 1];
    }
};
\end{lstlisting}

\solstep{Trace}
\begin{tracebox}[$R$ level sums, \figref{c03-running}]
Levels give $\text{sums}=[1,5,15]$. $k=2$: after a descending select, the array's index $1$
holds $5$ (order $\{15,5,1\}$) $\Rightarrow$ return $\mathbf{5}$. $k=3\Rightarrow 1$;
$k=4 > $ size $\Rightarrow -1$. \checkmark\ \probref{c24-kthlevel}.
\end{tracebox}

\solstep{Complexity}
\begin{complexitybox}
BFS $\bigO{n}$; \code{nth\_element} is $\bigO{\ell}$ average ($\ell$ levels), a full sort
$\bigO{\ell\log\ell}$ -- both dominated by $\bigO{n}$. Space $\bigO{n}$ (queue) $+\bigO{\ell}$.
\code{long long} prevents overflow on large-magnitude levels. Verdict: \ac.
\end{complexitybox}

\solstep{Pitfalls and Variants}
\begin{warningbox}
\code{nth\_element} with \code{greater<long long>} partitions so that \code{sums[k-1]} is the
$k$-th \emph{largest}; using the default (ascending) \code{nth\_element} then indexing
\code{[k-1]} yields the $k$-th \emph{smallest} -- the inverted-comparator mistake. Also guard
\code{sums.size()} $< k$ before touching \code{[k-1]}, and keep level sums in \code{long long}.
\end{warningbox}

\begin{insightbox}[Key takeaways -- three things to remember]
\begin{enumerate}[nosep, leftmargin=1.6em]
  \item Once you have the \emph{depth sequence}, ``$k$-th largest over levels'' is a plain
        array order-statistic -- the tree is already gone.
  \item $\ell \le n$ and often tiny, so sorting the level sums is free; select if you want the
        linear constant.
  \item The interesting bugs are the size guard ($k>\ell$) and the comparator direction.
\end{enumerate}
\end{insightbox}

\begin{hardeningbox}
\begin{itemize}[nosep, leftmargin=1.4em]
  \item \textbf{Stream / huge $\ell$}: if levels were too many to store, keep a size-$k$
        min-heap across the level sums -- the standard ``$k$-th largest of a stream'' variant.
  \item \textbf{DFS instead of BFS}: fill \code{level[d]} during a DFS (\coderef{c24-path-stack})
        and read the array -- same result, but $\bigO{h}$ stack on a skewed tree.
\end{itemize}
\end{hardeningbox}

\subsection{Practice Problems}

\begin{stuckbox}
\begin{itemize}[nosep, leftmargin=1.4em]
  \item Did you \emph{undo} the path aggregate on exit, or only push it? (Trap 1.)
  \item Is your reported level $1$-indexed when the problem wants it? (Trap 2,
        \probref{c24-levelsum}.)
  \item Did you guard ``fewer than $k$ levels $\to -1$'' before indexing? (\probref{c24-kthlevel}.)
  \item Are level sums accumulated in \code{long long} against negative/large values?
\end{itemize}
\end{stuckbox}

\begin{practiceset}
  \practice{Linked List in Binary Tree}{LeetCode 1367}
    {https://leetcode.com/problems/linked-list-in-binary-tree/}{LC Medium L3}
    {Does the \code{head} list appear as a downward \emph{root-to-node path prefix} in the
    tree? DFS carrying the list along the recursion stack: match node-by-node, backtrack
    (undo) when it fails -- the path-stack discipline of \secref{c24-path}. Start the list
    match at every tree node, or thread the \code{head} pointer down and allow a fresh start.}
\end{practiceset}

\begin{recallbox}[10-second refresher -- 5.3 Path \& Depth Sequence]
The DFS \emph{stack} is the root-to-node path: push the aggregate on entry, pop on exit
(never forget the pop). Bucket nodes by $\depth$ for the \emph{depth sequence} -- one array,
then scan/sort it (LC 1161 heaviest level, LC 2583 $k$-th largest level sum). The in/out
difference sweep reads root-path sums off the tour (\propref{c24-sweep}) -- the dual of
subtree flattening.
\selftest{
  \item Write the DFS that maintains a root-path sum with a matching undo on exit
        (\coderef{c24-path-stack}).
  \item State the in/out difference array and prove it returns the root-path sum
        (\propref{c24-sweep}).
  \item On $G$ with $w[v]=v$, list the level sums and name the heaviest 1-indexed level.
  \item Give the $2$-nd largest level sum of $R$ and say when the answer is $-1$.
  \item Explain why filling \code{level[d]} by DFS and sorting by $\depth$ are the same
        ``depth sequence''.
}
\end{recallbox}
